% ======================================================================
% 20-ssl.tex -- SPEC 0055, frames 16-23.
%
% §4.2 SSL association pretraining. Eight frames: divider; pretext;
% protocol; result; root cause; coordinate solvability; the frozen DynaCLR
% pivot; status and open branches. No SSL or contrastive tutorial.
% ======================================================================

% ----------------------------------------------------------------------
% Frame 16 -- divider.
% ----------------------------------------------------------------------
\talksectiondivider{\S4.2: Can warps manufacture association labels?}{One question: can a warp supply the label for free?}

% ----------------------------------------------------------------------
% Frame 17 -- why the pretext is not just another augmentation.
% (a) Why: separate the pretext from Trackastra's augmentation stack; the
%     claim is about the supervision source, not the warp type.
% (b) Tried: one frame, warp T, synthetic next frame, labels i -> T(i).
% (c) Novel: annotation-free association-label construction for the
%     low-label regime.
% (d) Interpretation: legitimate SSL idea; target is low-label. Threat: the
%     synthetic pair is not real motion.
% ----------------------------------------------------------------------
\begin{frame}
  \talktitle{Why the warp pretext is not just another augmentation}
  \begin{center}
  \begin{tikzpicture}[
    card/.style={draw=talkblue, fill=talkblue!6, rounded corners=2pt,
      text width=0.86\linewidth, align=center, inner sep=5pt, font=\small}
  ]
    \node[card] (cardone) {Trackastra augmentations perturb annotated windows};
    \node[card, below=4mm of cardone] (cardtwo) {Warp derives $i\to T(i)$\scite{wang2019learning,jabri2020spacetime}};
    \node[card, below=4mm of cardtwo] (cardthree) {No tracking annotations; targets the low-label regime};
  \end{tikzpicture}
  \end{center}
\end{frame}

% ----------------------------------------------------------------------
% Frame 18 -- the protocol.
% (a) Why: test the pretext under a pre-registered low-label rule, not post
%     hoc.
% (b) Tried: pretrain on trpL unlabeled frames (evaluation frames excluded),
%     fine-tune on five labeled pairs, random init and compute-matched
%     control, pre-registered success rule.
% (c) Novel: controlled low-label pilot with a compute-matched control and a
%     held-out condition.
% (d) Interpretation: if the rule fails, report failure; no protocol tweaks.
%     Threat: prototype scale.
% ----------------------------------------------------------------------
\begin{frame}
  \talktitle{What we tried: \texttt{trpL} held out, 5 labeled pairs, compute-matched control}
  \vfill
  \begin{center}
  \begin{tikzpicture}[
    card/.style={draw=talkblue, fill=talkblue!6, rounded corners=2pt,
      text width=0.40\linewidth, align=center, inner sep=5pt, font=\small}
  ]
    \node[card] (cardone) {\texttt{trpL} held out};
    \node[card, right=4mm of cardone] (cardtwo) {5 labeled pairs};
    \node[card, below=4mm of cardone] (cardthree) {random init + compute-matched control};
    \node[card, below=4mm of cardtwo] (cardfour) {pre-registered success rule, fixed before the run};
  \end{tikzpicture}
  \end{center}
  \vfill
\end{frame}

% ----------------------------------------------------------------------
% Frame 19 -- the result.
% (a) Why: does pretraining transfer to the held-out condition?
% (b) Tried: three arms, validation balanced accuracy.
% (c) Novel: the compute-matched control isolates the pretext from extra
%     compute.
% (d) Interpretation: all three arms 0.500; the rule was not met. Threat:
%     balanced accuracy at chance can hide calibration differences.
% ----------------------------------------------------------------------
\begin{frame}
  \talktitle{Result: all three arms 0.500}
  \begin{center}
  \small
  \begin{tabular}{lc}
    \toprule
    Arm & Validation balanced accuracy\\
    \midrule
    Random init & 0.500\\
    Warp pretext & 0.500\\
    Matched control & 0.500\\
    \bottomrule
  \end{tabular}
  \end{center}
  \vspace{4mm}

  \begin{center}
  {\footnotesize Chance is 0.500; the pre-registered success rule was not met.}
  \end{center}
\end{frame}

% ----------------------------------------------------------------------
% Frame 20 -- root cause.
% (a) Why: a null result is only useful with a root cause.
% (b) Tried: loss and entropy analysis, pool dependence, input audit,
%     true-identity accuracy.
% (c) Novel: root-cause evidence for objective-level collapse, not an
%     implementation bug.
% (d) Interpretation: loss 0.161 at the entropy floor; the same recipe
%     learns on another pool; input audit rules out an artifact. Threat:
%     one prototype configuration.
% ----------------------------------------------------------------------
\begin{frame}
  \talktitle{Why it collapsed: $\sim 1/N$ positives and an all-negative optimum}
  \begin{center}
  \begin{tikzpicture}[
    card/.style={draw=talkblue, fill=talkblue!6, rounded corners=2pt,
      text width=0.86\linewidth, align=center, inner sep=5pt, font=\small}
  ]
    \node[card] (cardone) {$\sim 1/N$ positives in the $N{\times}N$ BCE};
    \node[card, below=4mm of cardone] (cardtwo) {all-negative is the analytic optimum};
    \node[card, below=4mm of cardtwo] (cardthree) {loss 0.161 at the entropy floor; control pool 0.10 vs 0.188 floor};
  \end{tikzpicture}
  \end{center}
  \vspace{2mm}

  \begin{center}
  {\footnotesize Input audit: targets match the true correspondence; view-2 appearance is correctly warped; identical code learns on another pool.}
  \end{center}
\end{frame}

% ----------------------------------------------------------------------
% Frame 21 -- coordinate solvability.
% (a) Why: are the warp labels themselves too easy or positionally solvable?
% (b) Tried: coordinate-only probe on 92 frames, 8 experiments, 15 warp
%     draws per frame, seed 42.
% (c) Novel: quantifies coordinate recoverability of the warp labels.
% (d) Interpretation: a property of the labels, not a verdict on the
%     pretext's training value. Threat: applied configuration only.
% ----------------------------------------------------------------------
\begin{frame}
  \talktitle{Coordinate solvability: 0.982 of warp labels recoverable from coordinates}
  \begin{center}
  \small
  \begin{tabular}{@{}lccc@{}}
    \toprule
    Configuration & ICP & Warp-family & Oracle\\
    \midrule
    Applied: affine $10^\circ$, jitter 4, dropout 0.1 & 0.854 & 0.982 & 0.999\\
    Wider: affine $25^\circ$, jitter 8, dropout 0.1 & 0.629 & 0.900 & 0.933\\
    \addlinespace
    \multicolumn{4}{@{}l}{\footnotesize Random baseline $\approx 0.03$ (all columns)}\\
    \bottomrule
  \end{tabular}
  \end{center}
  \vspace{2mm}

  \begin{center}
  {\footnotesize This is a property of the labels, not a verdict on the pretext's training value.}
  \end{center}
\end{frame}

% ----------------------------------------------------------------------
% Frame 22 -- the pivot.
% (a) Why: after the pretext failure, build a legitimate SSL component
%     within scope.
% (b) Tried: DynaCLR-style NT-Xent ScaledCNN, 1.47M parameters, 128-D,
%     frozen; evaluated under Achard's protocol.
% (c) Novel: small from-scratch encoder versus large frozen foundation
%     models; feature-encoder pretraining, not association-head pretraining.
% (d) Interpretation: the pivot is a feature encoder, not the proposed
%     association-head pretraining. Threat: pairing uses tracking-derived
%     masks.
% ----------------------------------------------------------------------
\begin{frame}
  \talktitle{What we did instead: DynaCLR encoder 1.47M, frozen}
  \begin{center}
  {\small DynaCLR \scite{hirata2025dynacl}; Achard protocol \scite{achard2025thesis}.\par}
  \vspace{2mm}
  \begin{tikzpicture}[
    box/.style={draw=talkblue, fill=talkblue!6, rounded corners=2pt,
      align=center, inner sep=5pt, font=\small, minimum height=9mm},
    arrow/.style={-{Latex[length=2mm]}, thick, draw=talkblue}
  ]
    \node[box] (patches) {Cell patches};
    \node[box, right=8mm of patches] (encoder) {ScaledCNN + BatchNorm};
    \node[box, right=8mm of encoder] (embedding) {128-D embedding};
    \node[font=\scriptsize\color{talkgray}, below=1mm of encoder] {frozen};
    \draw[arrow] (patches) -- (encoder);
    \draw[arrow] (encoder) -- (embedding);
  \end{tikzpicture}
  \end{center}
  \vspace{3mm}

  \begin{center}
  {\footnotesize 1.47M parameters versus SAM2.1's 80.8M.}\par
  \vspace{1mm}
  {\footnotesize This is a feature encoder, not the proposed association-head pretraining.}
  \end{center}
\end{frame}

% ----------------------------------------------------------------------
% Frame 23 -- status and open branches.
% (a) Why: make the decision and the remaining work explicit.
% (b) Tried: one applied configuration; no full-scale run, no family sweep.
% (c) Novel: none claimed.
% (d) Interpretation: deprioritized as a scope decision, not disproven.
% ----------------------------------------------------------------------
\begin{frame}
  \talktitle{Deprioritized, not disproven; two branches remain open}
  \begin{columns}[t]
    \begin{column}{0.46\textwidth}
      \centering
      {\small\textbf{Decision}}\par\vspace{2mm}
      \begin{tikzpicture}[
        card/.style={draw=talkblue, fill=talkblue!6, rounded corners=2pt,
          text width=0.92\linewidth, align=center, inner sep=4pt, font=\small}
      ]
        \node[card] (one) {Deprioritized as a scope decision};
        \node[card, below=3mm of one] (two) {Not a measured rejection};
        \node[card, below=3mm of two] (three) {No full-scale run};
      \end{tikzpicture}
    \end{column}
    \begin{column}{0.46\textwidth}
      \centering
      {\small\textbf{Still open}}\par\vspace{2mm}
      \begin{tikzpicture}[
        card/.style={draw=talkorange, fill=talkorange!6, rounded corners=2pt,
          text width=0.92\linewidth, align=center, inner sep=4pt, font=\small}
      ]
        \node[card] (one) {Distortion families};
        \node[card, below=3mm of one] (two) {Low-label vs full-label sweep};
        \node[card, below=3mm of two] (three) {Only one applied configuration};
      \end{tikzpicture}
    \end{column}
  \end{columns}
\end{frame}
